\section*{Chapter \ref{ch:g12:early-earth} --- The Early Earth and Its Atmosphere}

\begin{solution}{exo:g12:early-earth:1}
(a) \omterm{def:g12:early-earth:primitive-atmosphere}{outgassing}; (b) a \omterm{def:g12:early-earth:stromatolite}{stromatolite}; (c) a \omterm{def:g12:early-earth:bif}{banded iron formation}; (d) \omterm{def:g12:early-earth:red-beds}{red
beds}; (e) the \omterm{def:g12:early-earth:goe}{Great Oxidation Event}.
\end{solution}

\begin{solution}{exo:g12:early-earth:2}
Zircons that show liquid water (4.4 billion years); the first \omterm{def:g12:early-earth:stromatolite}{stromatolites}
of the Pilbara (3.43); the Hamersley \omterm{def:g12:early-earth:bif}{banded iron formations} (more than
2.5); the end of the \omterm{def:g12:early-earth:goe}{Great Oxidation Event} (2.06); the end of most \omterm{def:g12:early-earth:bif}{banded iron
formations} (1.8).
\end{solution}

\begin{solution}{exo:g12:early-earth:3}
At $t = 0$, $L = L_\odot/(1 + 0.4) = L_\odot/1.4 \approx 0.714\,L_\odot$:
about \qty{29}{\percent} dimmer, nearly a third less bright than today.
\end{solution}

\begin{solution}{exo:g12:early-earth:4}
That the air then held almost no oxygen. Pyrite is destroyed by oxygen;
its grains could be rolled along by rivers and deposited intact only if the
air and the water at the surface lacked it.
\end{solution}

\begin{solution}{exo:g12:early-earth:5}
Water vapour, which condensed into the ocean, and carbon dioxide, which
stayed in the air in much larger amounts than today and warmed the Earth by
its greenhouse effect.
\end{solution}

\begin{solution}{exo:g12:early-earth:6}
3 billion years ago: $1 - t/t_\odot = 3/4.567 \approx 0.657$, $L \approx
1/1.263 \approx 0.792\,L_\odot$, $T \approx 255 \times 0.792^{1/4} \approx
\qty{241}{K}$. 2.45 billion years ago: $1 - t/t_\odot \approx 0.536$, $L
\approx 1/1.215 \approx 0.823\,L_\odot$, $T \approx 255 \times 0.823^{1/4}
\approx \qty{243}{K}$.
\end{solution}

\begin{solution}{exo:g12:early-earth:7}
\qty{1}{kg} of iron is $1000/55.85 \approx 17.9$ moles; one molecule of
oxygen takes four ions of iron, so $17.9/4 \approx 4.48$ moles of
\ce{O2}, that is $4.48 \times 32.00 \approx \qty{143}{g}$ of oxygen.
\end{solution}

\begin{solution}{exo:g12:early-earth:8}
(1) $2.46 - 2.06 = 0.40$ billion years, 400 million years. (2) From 4.54
to 2.46 billion years ago: $2.08/4.54 \approx \qty{46}{\percent}$ of the
Earth's history. (3) About $3.43 - 2.46 \approx 1$ billion years: the oxygen
the microbes made was used up for a long time before it could stay in the
air.
\end{solution}

\begin{solution}{exo:g12:early-earth:9}
The \omterm{def:g12:early-earth:bif}{banded iron formation} shows a deep sea without oxygen, which could carry
dissolved iron, under a shallow sunlit sea with enough oxygen to precipitate
it. The red bed, laid down on land, shows oxygen in the air and in the water
at the surface, enough to rust the iron of its grains.
\end{solution}

\begin{solution}{exo:g12:early-earth:10}
In Hamelin Pool, Shark Bay, Western Australia, whose water is twice as
salty as usual sea water. A mat of microbes, mainly cyanobacteria, traps and
binds grains of sediment and makes carbonate grow; a new mat grows on top,
and so on, layer upon layer.
\end{solution}

\begin{solution}{exo:g12:early-earth:11}
(1) $21 \times 10^{-5} = \qty{0.00021}{\percent}$. (2) $21/10 =
\qty{2.1}{\percent}$.
\end{solution}

\begin{solution}{exo:g12:early-earth:12}
The \omterm{def:g10:habitable-planet:equilibrium-temperature}{equilibrium temperature} is unchanged if $L(1 - A')$ equals $1 - A$:
$1 - A' = 0.706/0.741 \approx 0.953$, so $A' \approx 0.047$. The Earth would
have had to reflect less than \qty{5}{\percent} of the sunlight, almost like a
black body, which an ocean planet with clouds and some ice cannot do. The
other way out is a stronger greenhouse effect.
\end{solution}

\begin{solution}{exo:g12:early-earth:13}
(1) $273 - 237 = \qty{36}{K}$ to reach freezing; $288 - 237 = \qty{51}{K}$ to
reach today's mean. (2) Today's greenhouse effect adds $288 - 255 =
\qty{33}{K}$: the young Earth needed a stronger one. (3) \qty{0.1}{bar} is a
tenth of the pressure of today's whole atmosphere, for carbon dioxide alone.
\end{solution}

\begin{solution}{exo:g12:early-earth:14}
(1) Before 2450 million years ago the signature exists, so the air held less
than $10^{-5}$ of today's oxygen; after 2090 million years it is absent, so
the air held more. (2) Yes: the \omterm{def:g12:early-earth:goe}{Great Oxidation Event} runs from about 2460 to
2060 million years ago, the same interval.
\end{solution}

\begin{solution}{exo:g12:early-earth:15}
(1) About $4567 - 60 = 4507$ million years ago. (2) At most $4507 - 4404 =
103$ million years after the impact: the Earth had a crust and water within
about a hundred million years of the impact.
\end{solution}

\begin{solution}{pb:g12:early-earth:1}
\textbf{1.} $4.567 - 4.0 = 0.567$ billion years.

\textbf{2.} $1 - 0.567/4.567 \approx 1 - 0.124 = 0.876$.

\textbf{3.} $L = 1/(1 + 0.4 \times 0.876) = 1/1.350 \approx 0.741$ of
today's.

\textbf{4.} $0.741 \times 1361 \approx \qty{1008}{W/m^2}$.

\textbf{5.} $1 - 0.741 \approx 0.26$: \qty{26}{\percent} dimmer.

\textbf{6.} $T_{\mathrm{eq}} = \left(S(1-A)/(4\sigma)\right)^{1/4}$.

\textbf{7.} With $A$ and $\sigma$ fixed, $T_{\mathrm{eq}}$ is a constant times
$S^{1/4}$, so $T'/T = (S'/S)^{1/4}$.

\textbf{8.} $0.741^{1/4} = \sqrt{\sqrt{0.741}} \approx \sqrt{0.861} \approx
0.928$.

\textbf{9.} $255 \times 0.928 \approx \qty{237}{K}$.

\textbf{10.} $237 - 273 \approx \qty{-36}{\celsius}$.

\textbf{11.} That there was liquid water at the surface, an ocean, 4.4
billion years ago.

\textbf{12.} In shallow water, as living \omterm{def:g12:early-earth:stromatolite}{stromatolites} do today.

\textbf{13.} The \omterm{def:g12:early-earth:bif}{banded iron formations} of the Hamersley Ranges, laid down
on a sea floor.

\textbf{14.} Yes: \qty{237}{K} is far below \qty{273}{K}; with today's
atmosphere the mean surface would have been frozen.

\textbf{15.} The rocks show liquid water under a Sun too faint to keep it
liquid with today's atmosphere: the \omterm{def:g12:early-earth:faint-young-sun}{faint young Sun problem}.

\textbf{16.} $288 - 255 = \qty{33}{K}$.

\textbf{17.} At least $273 - 237 = \qty{36}{K}$, more than today's.

\textbf{18.} $A' = 1 - 0.706/0.741 \approx 0.047$: an almost black planet,
not plausible.

\textbf{19.} Carbon dioxide, much more abundant than today; methane, made by
early microbes, may have helped.

\textbf{20.} The Earth's \omterm{def:g10:habitable-planet:equilibrium-temperature}{equilibrium temperature} four billion years ago was
\textbf{about 237 K}, some \qty{36}{\celsius} below freezing.
\end{solution}
